\section{Shallow next-to-leading coefficients}
\label{sec:shallow-top-coefficient}

\begin{corollary}\label{cor:shallow-top-split}
Let \(p\) be any prime, \(n\ge2\), and \(O=\Z_p\) or
\(\Z_{(p)}\).  If
\[
 F=t^n+pb\,t^{n-1}+\sum_{i=0}^{n-2}a_it^i,
 \qquad b\in O,\quad a_i\in p^2O,
\]
the generic kernel for \(O[t]/F\) contains a countable direct sum
of exact order-\(n\) classes which splits from the whole
\(K_0(\Lat_{O[t]/F}(S_p))\).  The retraction annihilates the
Young span.  The analogous cyclic-group statement also holds.
\end{corollary}
\begin{proof}
If \(b\in pO\), then \(r=\infty\); otherwise \(r=n-1\).
In either case \(N=n\), because \(p(n-1)\ge n\).
Apply Theorem~\ref{thm:singular-distinguished-monic}.
\end{proof}

For odd primes the following direct calculation describes another
coordinate system and an explicit cubic companion lattice.
Its particular preimages are not identified with the mesh section
in the preceding corollary.

\begin{theorem}\label{thm:shallow-top-coefficient}
Let $p$ be odd, $n\ge2$, $O=\Z_p$ or $\Z_{(p)}$, and
\[
 F(t)=t^n+pb\,t^{n-1}+\sum_{i=0}^{n-2}a_it^i,
 \qquad b\in O,\quad a_i\in p^2O.
\]
Put $A=O[t]/F(t)$.  There is a homomorphism
\[
 K_0(\Lat_A(S_p))\longrightarrow\bigoplus_{j\ge1}\Z/n
\]
which annihilates the degree-$p$ Young permutation span and whose
restriction to the kernel of generic specialization is surjective.
The analogous assertion holds for $C_p$.
In particular, $\Phi_A$ is not surjective in degree $p$.
No annihilation relation or splitting of the preimages is asserted.
\end{theorem}

The proof changes the coefficient test modules, so that every test
map and its suspension still lift.  It does not apply the binomial
Bockstein to unavailable truncated-polynomial lifts.

\begin{proof}
If $n=2$, replacing $t$ by $u-pb/2$ identifies $A$ with
$O[u]/(u^2-d)$, where $d=p^2b^2/4-a_0\in p^2O$.
Theorem~\ref{thm:binomial-depth-two-all-exponents} proves the result.
Hence assume $n\ge3$, initially over $O=\Z_p$.
Write $k=\F_p$, $W=O/p^2O$, $R=k[t]/t^n$, and
$\widetilde R=A/p^2A$.

\paragraph{A palindromic chain of coefficient lifts.}
Choose monic linear polynomials $\ell_1,\ldots,\ell_n\in W[t]$.
For $n=2m+1$, let all of them be $t$, except
$\ell_{m+1}=t+pb$.
For $n=2m$, let all of them be $t$, except
$\ell_m=\ell_{m+1}=t+pb/2$.
Their sequence is palindromic, and
\[
 \prod_{r=1}^n\ell_r=t^n+pb\,t^{n-1}=F\pmod{p^2}.
\]
For $0\le i\le n$ put $f_i=\prod_{r=1}^i\ell_r$, with
$f_0=1$.  Then
\begin{equation}\label{eq:shallow-palindrome}
 f_i\mid f_j\quad(i\le j),\qquad
 f_i f_{n-i}=f_n,
 \qquad \widetilde R=W[t]/f_n.
\end{equation}
Let $X_i=R/(t^i)$ and
$\widetilde X_i=W[t]/f_i$, for $1\le i\le n$.
These are $W$-free lifts of the $X_i$.

Every $R$-map $X_i\to X_j$ lifts to a
$\widetilde R$-map.  If $i\le j$, its value at $1$ has the
form $t^{j-i}h(t)$; lift it to
$(f_j/f_i)\widetilde h(t)\bmod f_j$, lifting each coefficient
of $h$ arbitrarily to $W$.
If $i\ge j$, lift the arbitrary image $h(t)$ of $1$ to
$\widetilde h(t)\bmod f_j$; the divisibility $f_j\mid f_i$
makes this well-defined.
Matrices of these maps cover arbitrary finite direct sums.
An isomorphism lifts to an isomorphism, since its reduction
is invertible between $W$-free modules of the same rank.

The coefficient syzygy also lifts exactly:
\begin{equation}\label{eq:shallow-lifted-syzygy}
 0\longrightarrow\widetilde X_i
 \xrightarrow{\ f_{n-i}\ }\widetilde R
 \longrightarrow\widetilde X_{n-i}\longrightarrow0.
\end{equation}
For injectivity, if $f_{n-i}a\in(f_i f_{n-i})$ in $W[t]$,
cancel the monic polynomial $f_{n-i}$ to get $a\in(f_i)$.
The cokernel is visibly $W[t]/f_{n-i}$.
The sequence is split over $W$ because the quotient is
$W$-free.  Its reduction is
$0\to X_i\xrightarrow{t^{n-i}}R\to X_{n-i}\to0$.

\paragraph{The full-category Bockstein and its suspension.}
For an $A$-free cyclic lattice $E$, use curvature
$U_E=(\widetilde T^p-I)/F\bmod F$.
The stable naturality, integral primary extraction, and
cohomological triviality of every component $E_f$, $f\ne z$,
in Lemma~\ref{lem:binomial-all-primary} remain valid.
Indeed, their proofs use trace telescoping and cancellation
of the monic $F$, with $F\bmod p=t^n$; no congruence modulo
$p^2$ enters those arguments.
For a selected component $E$, its exact residual Tate
complex has free $R$-terms.  Thus
Lemma~\ref{lem:binomial-all-representability} supplies
\[
 M_E=\operatorname{coker}(T-I)^*,\qquad
 \underline{\operatorname{Hom}}_R(M_E,X)
 \cong\widehat H^0(C_p,(E/pE)\otimes_R X),
\]
and the contravariant short exact sequence of transpose
cokernels associated to every admissible sequence.

For a direct sum of the chosen test lifts, use
\[
 0\longrightarrow(E/pE)\otimes_R X
 \xrightarrow{p}(E/p^2E)\otimes_{\widetilde R}\widetilde X
 \longrightarrow(E/pE)\otimes_R X\longrightarrow0.
\]
It is exact because $E/p^2E$ is $\widetilde R$-free and
$\widetilde X$ is $W$-free.  Its Tate boundary is $\beta_X$.
The preceding lift of every map gives naturality on the
entire finite $R$-module category.  No functorial choice of
map lifts is needed: every such lift gives the naturality
square with the same prescribed maps on its two reductions.
A change of cyclic decomposition lifts to an isomorphism,
so the resulting operator is independent of that presentation.

Use the cyclic homotopy
$h_{\rm even}=I$, $h_{\rm odd}=-Q_p(T)$, with
$Q_p(z)=(N(z)-p)/(z-1)$.
The identity $dh+hd=pI$ gives $\beta h+h\beta=I$.
Since $p$ is odd, $h^2=-Q_p(1)=0$ on Tate cohomology, and
$P=\beta h$ is idempotent.

We check the boundary signs with the new lifts explicitly.
For an admissible lattice sequence, or for
\eqref{eq:shallow-lifted-syzygy} after tensoring with $E/p^2E$,
choose a degreewise $W$-linear splitting of the complexes.
Write the off-diagonal differential block as $Z$.
Then $d'Z+Zd''=0$ over $W$.
For a lifted cycle with $d''z=pw$, we have
$d'(Zz)=-pZw$; dividing by $p$ shows
$\beta'\partial+\partial\beta''=0$.
Writing the off-diagonal block of $h$ as $Y$, the off-diagonal
part of $dh+hd=pI$ is
\[
 d'Y+Yd''+h'Z+Zh''=0.
\]
It gives $h'\partial+\partial h''=0$ on cohomology.
Consequently $P'\partial=\partial P''$ in both kinds of
sequence.  The middle coefficient test $R$ has zero Tate
cohomology, so the reduced syzygy boundaries identify the
two parities with stable suspension.

Yoneda therefore gives a stable idempotent $P_E$ of $M_E$
which commutes with all three arrows of the transpose
triangle, including its suspension arrow.
Let $M_E^+$ be its image in the idempotent-complete stable
category of $R$.  The stable-length and direct-summand
arguments of Section~\ref{sec:binomial-depth-two-all} now
prove that
\[
 \eta_f(L)=-\dim_k M_{L_f}^+\pmod n\qquad(f\ne z)
\]
is additive on every admissible sequence.
These arguments use the Frobenius stable-category theorem
and the direct-summand theorem for triangles in exactly the
scopes already established there
\citep[Section~13.4, Theorem~13.7, author p.~50]{Buhler2010}
\citep[Lemma~1.6 and proof, pp.~821--822]{BalmerSchlichting2001}.
Choice independence holds on the full category: maps
factoring through $A[C_p]$ act as zero on every test functor,
and the natural $\beta,h$ intertwine the stable isomorphisms
between choices of curvature-primary components.
Each class has finite primary support.  Constant $O$-lattices
have zero curvature and hence all coordinates zero.

\paragraph{Witnesses and their actual selected stable module.}
Write $F=a_0+tH(t)$, and work over a finite unramified
extension $O_s/O$.  If $a_0=p^v u\ne0$, take
\[
 \alpha=1+p^{v-1}y,\qquad c=(\alpha^p-1)/a_0,
\]
where $y$ is a unit.  Then $c$ is a unit and
$\overline c=\overline y/\overline u$ can be prescribed
arbitrarily in $k_s^\times$.
If $a_0=0$, take $\alpha=1$ and any unit $c$.
In both cases $\alpha^p-ca_0=1$.
Let $C_\alpha$ be the companion of
$q_\alpha(z)=(z^p-\alpha^p)/(z-\alpha)$,
with last coordinate row $a$ and first coordinate column $b_0$.
The matrix
\[
 T=\begin{pmatrix}\alpha&ct\,a\\H(t)b_0&C_\alpha\end{pmatrix}
\]
has characteristic polynomial $z^p-1-cF(t)$ before imposing
$F=0$, so $\widetilde T^p-I=cF(t)I$.
It defines a lattice $L$ of scalar curvature $c$.

Here $H\bmod p=t^{n-1}$.
Present the residual transpose cokernel by
$x,y_1,\ldots,y_{p-1}$.
If $m=p-1$ and
$\overline q_\alpha=z^m+q_{m-1}z^{m-1}+\cdots+q_0=(z-1)^m$,
the top row gives $\overline c\,t y_m=0$; the subsequent
companion rows give
$y_{r-1}-y_r-q_{r-1}y_m=0$ for $2\le r\le m$.
Eliminate $y_1,\ldots,y_{m-1}$.
The remaining relation is
\[
 t^{n-1}x-(1+q_{m-1}+\cdots+q_0)y_m
 =t^{n-1}x-\overline q_\alpha(1)y_m=t^{n-1}x=0.
\]
Thus, over $R_s=k_s[t]/t^n$,
\begin{equation}\label{eq:shallow-witness-transpose}
 M_L\cong R_s/(t^{n-1})\,x\ \oplus\ R_s/(t)\,y_m.
\end{equation}

Since $n\ge3$, the palindromic chain has $f_1=t$.
Consequently the chosen lift of the residue test $k_s$ has
$t=0$ even modulo $p^2$.
The upper off-diagonal block vanishes there, as does the
lower one modulo $p^2$, because $H(0)=a_1\in p^2O$.
On the trivial even Tate class, $h=I$ and
$\beta_{\rm odd}=N(\alpha)/p\equiv1\pmod p$.
The last congruence follows by expanding $N(1+w)$ with
$w\in pO_s$ and dividing by $p$.
The even cyclotomic class is represented by
$(C_\alpha-I)^{p-2}b_0$ modulo $p$; since $p>2$, its image
under $h_{\rm even}=I$ belongs to
$\operatorname{im}(C_\alpha-I)$ and is zero in odd cohomology.
Thus $P$ has values $1,0$ on the two residue-test classes.
The idempotent-lifting and Krull--Schmidt argument applied
to \eqref{eq:shallow-witness-transpose} identifies its image
as $R_s/(t^{n-1})$ up to stable isomorphism.
Therefore the negative-length coordinate of $L$ is
$-s(n-1)=s\pmod n$.

The witness is generically regular without a reducedness
assumption on the generic algebra.  Put $B=A_s[1/p]$.
The $B$-free representation is projective over $B[C_p]$
by averaging.  Over every residue field of $B$, its
characteristic polynomial $z^p-1$ is separable and gives
the regular representation.
The commutative Artinian algebra $B[C_p]$ is finite etale
over $B$.  The projective module has rank one on each
of its local factors, so is free of rank one over $B[C_p]$.
Thus $[L]-[A_s[C_p]]$ vanishes generically.

Choose distinct prime degrees $s_j>n(p-1)$, and residue
parameters of those degrees.  The scalar orbits of these
parameters under $\F_p^\times$ are free: equality between
$a\overline c$ and a Frobenius conjugate of $\overline c$
would give $a^{s_j}=1$ after taking the norm, hence $a=1$.
Different degrees give disjoint orbits.
Restrict coefficients to $A$ and induce to $S_p$.
Changing the cyclic generator to its $a$th power multiplies
the stable curvature by $a$; Mackey restriction supplies
each normalizer twist once, and all remaining terms are
regular cyclic lattices.
Accordingly, the $f_j$-coordinates of the generic-zero
classes $[\operatorname{Ind}L_j]-s_j[A[S_p]]$ are diagonal,
with value $s_j\pmod n$ on the diagonal.
Normalize by $s_j^{-1}$.  Finite primary support gives the
claimed map to the direct sum, and these witnesses prove
surjectivity from its generic kernel.
Every Young lattice and its cyclic restriction is constant,
so the Young span is annihilated.

For $O=\Z_{(p)}$, compose the completed coordinates with
completion.  Lift the separable residue parameter polynomial
to obtain a finite free local etale $O$-algebra $O_s$.
The finite formulas for $\alpha,c$ are integral over $O_s$
and define the same companions there.
The preceding Artinian generic argument proves their generic
vanishing before completion, and their completed coordinate
values are unchanged.  They therefore give the desired
localized preimages.  No mesh relation is used.
\end{proof}

\begin{example}\label{ex:shallow-wild-cubic}
At $p=3$, put $A=\Z_3[t]/(t^3+3t^2+9)$.
Its Newton polygon has a single segment of slope $-2/3$.
The coprime-denominator argument proves irreducibility and
total ramification of degree three
\citep[Corollary~3.18, p.~49]{WilliamsWuthrich2025}.
Here is also an explicit normalization calculation.
Set $\pi=t^2/3$.  The defining equation gives
\[
 \pi(t+3)=-3,\qquad
 \pi^3=3(\pi+1)^2,
 \qquad t=-\pi^2+3\pi+3.
\]
The polynomial $\pi^3-3\pi^2-6\pi-3$ is Eisenstein;
thus the normalization is $S=\Z_3[\pi]$
\citep[Theorem~5.9, pp.~67--68]{WilliamsWuthrich2025}.
Since $t^2=3\pi$, the two bases satisfy
\[
 A=\Z_3\oplus3\Z_3\pi\oplus\Z_3\pi^2\subset S.
\]
The index is three, so this is a proper local order in a
totally ramified cubic field.

Choose $\alpha=4$, $c=7$, so $\alpha^3-9c=1$.
One explicit witness is
\[
 T=\begin{pmatrix}
 4&0&7t\\t^2+3t&0&-16\\0&1&-4
 \end{pmatrix},\qquad T^3=I.
\]
Its curvature residue is $1$, with free scalar orbit under
$\F_3^\times$.
The induced class $[\operatorname{Ind}_{C_3}^{S_3}L]-[A[S_3]]$
vanishes generically and has coordinate $1\pmod3$.
This proves it is nonzero; its exact order is not asserted.
The countable family is obtained with the unramified
parameters in the proof.
\end{example}
